\documentclass[UTF8,a4paper,11pt]{ctexart}

\usepackage{geometry}
\usepackage{amsmath,amssymb,bm}
\usepackage{graphicx}
\usepackage{booktabs}
\usepackage{caption}
\usepackage{subcaption}
\usepackage{float}
\usepackage{xcolor}
\usepackage{hyperref}

\geometry{left=2.4cm,right=2.4cm,top=2.4cm,bottom=2.4cm}
\hypersetup{colorlinks=true,linkcolor=blue,urlcolor=blue,citecolor=blue}
\setlength{\parskip}{0.35em}
\setlength{\parindent}{2em}

\newcommand{\figfile}[1]{\detokenize{#1}}
\newcommand{\caseplot}[2]{
  \begin{figure}[H]
    \centering
    \begin{subfigure}{0.32\linewidth}
      \centering
      \includegraphics[width=\linewidth]{\figfile{../runs/ppt_formula_exact_only/#1/roe_cfl_0p5_nx_501/plots/rho.png}}
      \caption{density}
    \end{subfigure}
    \begin{subfigure}{0.32\linewidth}
      \centering
      \includegraphics[width=\linewidth]{\figfile{../runs/ppt_formula_exact_only/#1/roe_cfl_0p5_nx_501/plots/u.png}}
      \caption{velocity}
    \end{subfigure}
    \begin{subfigure}{0.32\linewidth}
      \centering
      \includegraphics[width=\linewidth]{\figfile{../runs/ppt_formula_exact_only/#1/roe_cfl_0p5_nx_501/plots/p.png}}
      \caption{pressure}
    \end{subfigure}
    \caption{#2: Roe 数值解与精确 Riemann 解对比。}
  \end{figure}
}

\title{Project 3 结果报告：Roe 格式求解一维 Euler 方程 Riemann 问题}
\author{}
\date{}

\begin{document}
\maketitle

\section{目标}

本报告只讨论课件 Roe 原式的一维 Euler 方程 Riemann 问题数值结果。核心目标是将 Roe finite-volume 数值解与精确 Riemann 解对比，检查典型激波管、膨胀波、接触间断和双激波问题中的波形、误差和失稳情况。

\section{数值方法}

一维 Euler 方程写成守恒形式：
\begin{equation}
    \frac{\partial \bm Q}{\partial t}
    +
    \frac{\partial \bm F(\bm Q)}{\partial x}
    =0,
    \qquad
    \bm Q=
    \begin{bmatrix}
        \rho\\ \rho u\\ \rho E
    \end{bmatrix},
    \quad
    \bm F=
    \begin{bmatrix}
        \rho u\\ \rho u^2+p\\ u(\rho E+p)
    \end{bmatrix}.
\end{equation}

离散格式采用 cell-centered finite volume 守恒更新：
\begin{equation}
    \bm Q_j^{n+1}
    =
    \bm Q_j^n
    -
    \frac{\Delta t}{\Delta x}
    \left(
        \hat{\bm F}_{j+1/2}
        -
        \hat{\bm F}_{j-1/2}
    \right).
\end{equation}

界面数值通量按课件 Roe 公式实现：
\begin{equation}
    \hat{\bm F}_{j+1/2}
    =
    \frac12
    \left[
        \bm F(\bm Q_{L,j+1/2})
        +
        \bm F(\bm Q_{R,j+1/2})
    \right]
    -
    \frac12
    |\tilde{\bm A}|_{j+1/2}
    \left(
        \bm Q_{R,j+1/2}
        -
        \bm Q_{L,j+1/2}
    \right).
\end{equation}

当前程序是一阶格式，界面左右状态直接取相邻控制体：
\begin{equation}
    \bm Q_{L,j+1/2}=\bm Q_j,
    \qquad
    \bm Q_{R,j+1/2}=\bm Q_{j+1}.
\end{equation}
程序显式构造 Roe 平均量和课件中的 \(\alpha_1,\ldots,\alpha_5\)，并按
\begin{equation}
    |\tilde{\bm A}|\Delta \bm Q
    =
    \begin{bmatrix}
        \alpha_4\\
        \tilde u\alpha_4+\alpha_5\\
        \tilde H\alpha_4+\tilde u\alpha_5
        -\tilde c^2\alpha_1/(\gamma-1)
    \end{bmatrix}
\end{equation}
组装耗散项。

\section{基本设置与后处理}

主线计算参数为：
\begin{itemize}
    \item 计算域：\([0,1]\)，初始间断位置 \(x_0=0.5\)；
    \item 网格：\(N_x=501\) 个 cell-centered 控制体；
    \item 气体参数：\(\gamma=1.4\)；
    \item CFL：0.5；
    \item 对比基准：已有 exact Riemann solver 生成的精确解。
\end{itemize}

后处理脚本读取 numerical 和 exact 的 Tecplot ASCII 文件。误差统计时先将 exact 解插值到 numerical cell-center 网格，再计算
\begin{equation}
    L_1=\int |q_{\mathrm{num}}-q_{\mathrm{exact}}|\,dx,
    \quad
    L_2=\left[\int (q_{\mathrm{num}}-q_{\mathrm{exact}})^2\,dx\right]^{1/2},
    \quad
    L_\infty=\max |q_{\mathrm{num}}-q_{\mathrm{exact}}|.
\end{equation}

\section{七组算例结果}

表 \ref{tab:baseline} 给出 Roe 原式与精确解的主线误差统计。case 04 在推进过程中产生非物理压力，因此没有最终 numerical/exact 曲线；该问题在讨论部分单独说明。

\begin{table}[H]
\centering
\caption{Roe 原式与精确解对比，\(N_x=501\), CFL=0.5。}
\label{tab:baseline}
\resizebox{\textwidth}{!}{%
\begin{tabular}{clrrrrrr}
\toprule
case & 名称 & \(\rho L_1\) & \(u L_1\) & \(p L_1\) & \(\rho L_\infty\) & \(u L_\infty\) & \(p L_\infty\)\\
\midrule
01 & Sod & \(5.897{\times}10^{-3}\) & \(7.402{\times}10^{-3}\) & \(4.441{\times}10^{-3}\) & \(8.175{\times}10^{-2}\) & \(3.282{\times}10^{-1}\) & \(5.827{\times}10^{-2}\)\\
02 & Lax & \(1.828{\times}10^{-2}\) & \(1.321{\times}10^{-2}\) & \(1.575{\times}10^{-2}\) & \(4.440{\times}10^{-1}\) & \(7.969{\times}10^{-1}\) & \(7.085{\times}10^{-1}\)\\
03 & Subsonic double expansion & \(7.375{\times}10^{-3}\) & \(2.422{\times}10^{-2}\) & \(2.452{\times}10^{-2}\) & \(1.114{\times}10^{-1}\) & \(1.011{\times}10^{-1}\) & \(8.145{\times}10^{-2}\)\\
04 & Sjogreen supersonic expansion & \multicolumn{6}{c}{unstable: nonpositive pressure}\\
05 & Contact double expansion & \(5.806{\times}10^{-3}\) & \(6.326{\times}10^{-3}\) & \(3.516{\times}10^{-3}\) & \(1.540{\times}10^{-1}\) & \(5.390{\times}10^{-2}\) & \(3.066{\times}10^{-2}\)\\
06 & Contact double shock & \(6.002{\times}10^{-3}\) & \(3.788{\times}10^{-3}\) & \(4.172{\times}10^{-3}\) & \(3.861{\times}10^{-1}\) & \(2.027{\times}10^{-1}\) & \(2.421{\times}10^{-1}\)\\
07 & Pure contact & \(1.261{\times}10^{-1}\) & \(6.150{\times}10^{-10}\) & \(3.010{\times}10^{-9}\) & \(4.142{\times}10^{0}\) & \(1.522{\times}10^{-9}\) & \(7.451{\times}10^{-9}\)\\
\bottomrule
\end{tabular}}
\end{table}

\begin{figure}[H]
    \centering
    \begin{subfigure}{0.48\linewidth}
        \centering
        \includegraphics[width=\linewidth]{\figfile{../runs/ppt_formula_exact_only/report_figures/baseline_l1_errors.png}}
        \caption{\(L_1\) 误差}
    \end{subfigure}
    \begin{subfigure}{0.48\linewidth}
        \centering
        \includegraphics[width=\linewidth]{\figfile{../runs/ppt_formula_exact_only/report_figures/baseline_linf_errors.png}}
        \caption{\(L_\infty\) 误差}
    \end{subfigure}
    \caption{Roe 原式与精确解的误差总览。case 04 因失稳未进入误差图。}
    \label{fig:error_overview}
\end{figure}

\caseplot{case_01}{case 01 Sod}
\caseplot{case_02}{case 02 Lax}
\caseplot{case_03}{case 03 Subsonic double expansion}
\caseplot{case_05}{case 05 Contact double expansion}
\caseplot{case_06}{case 06 Contact double shock}
\caseplot{case_07}{case 07 Pure contact}

\section{结果讨论}

整体来看，Roe 原式在 case 01、02、03、05、06 中能够捕捉主要波系结构，包括激波、膨胀波和接触间断。由于当前实现是一阶 Godunov 型格式，间断附近不可避免有数值扩散，因此 \(L_\infty\) 误差明显大于 \(L_1\) 误差。这个现象在 Lax 和 Contact shock 等含强间断的算例中尤为明显。

case 07 是纯接触间断。速度和压力在精确解中几乎保持常值，因此程序给出的 \(u\) 和 \(p\) 误差接近机器精度；但密度间断被一阶格式明显扩散，导致 \(\rho L_1\) 和 \(\rho L_\infty\) 较大。这说明 Roe 格式能保持接触波的速度和压力关系，但一阶重构仍会扩散密度跳跃。

case 04 在推进中出现非物理压力。这不是后处理错误，而是 Roe 原式对强膨胀、近真空问题不满足 positivity preservation 的典型表现。程序没有强行截断密度或压力，因此这个失败应作为方法局限写入报告。

\section{结论}

本实验按课件 Roe 公式显式实现了一阶 cell-centered FVM，并完成了七组 Riemann 问题与精确解的对比。结果表明，Roe 格式能够较好捕捉多数典型波系；但在一阶格式下，激波和接触间断附近存在明显数值扩散，\(L_\infty\) 误差主要由间断位置附近贡献。

case 04 的强膨胀问题暴露了 Roe 原式不保正性的缺陷。若后续课程允许进一步改进，应优先考虑保正的近似 Riemann solver 或 limiter，而不是在本报告中额外加入课件未要求的修正项。

\end{document}
